\section{X-Rays and Widths Have the Same Stability Rate}

\begin{theorem}[X-ray and width stability rates are equal]
\label{thm:final}
For convex bodies with support functions in $C^{1,\alpha}(S^{D-1})$ and $m$ measurement directions:
\[
  E_m^*(\mathrm{X\text{-}rays}) = \Theta(E_m^*(\mathrm{widths})) = \Theta(M \cdot m^{-(1+\alpha)/(D-1)}).
\]
\end{theorem}

\begin{proof}
\textbf{Upper bound.}

Near each measurement direction $u_i$: the X-ray data constrains $h_K$ (through the chord-length-to-boundary relationship and convexity). The constraint is stable when chords are long (directions close to $u_i$).

Between measurement directions: there are angular gaps of width $\omega_m \leq C_D m^{-1/(D-1)}$ (covering radius on $S^{D-1}$). In these gaps, no measurement of any kind provides direct information. The reconstruction of $h_K$ in the gap relies entirely on the regularity constraint $h_K \in C^{1,\alpha}$. The interpolation error across a gap of angular width $\omega$ is $O(M \omega^{1+\alpha})$ (this is the definition of $C^{1,\alpha}$: the function can deviate by at most $M\omega^{1+\alpha}$ from its value at the gap boundary, where $M = \|h_K\|_{C^{1,\alpha}}$).

Therefore: $\delta_H(K, L) \leq C M \omega_m^{1+\alpha} = C M m^{-(1+\alpha)/(D-1)}$.

This bound applies to both X-rays and widths (both have the same angular gaps).

\textbf{Lower bound.}

By the optimal recovery / Kolmogorov $n$-width theorem: for any $m$ measurement directions, there exists $g \in C^{1,\alpha}(S^{D-1})$ with $\|g\|_{C^{1,\alpha}} \leq M$, $g(u_i) = g(-u_i) = 0$ for all $i$, and $\|g\|_\infty \geq c M m^{-(1+\alpha)/(D-1)}$.

For widths: $w_K(u_i) = w_{K+g}(u_i)$ since $g(u_i) + g(-u_i) = 0$. So $K$ and $K+g$ are width-indistinguishable. Lower bound: $c M m^{-(1+\alpha)/(D-1)}$.

For X-rays: the chord-length perturbation at direction $u_i$ and offset $x$ depends on $g$ near $u_i$ (not just at $u_i$). Since $g(u_i) = 0$: $|g(v)| \leq \|\nabla g\|_\infty \cdot |v - u_i|$ for $v$ near $u_i$. The chord perturbation: $|\delta C(u_i, x)| \leq C \|\nabla g\|_\infty \cdot \kappa R$.

Now: $g$ is a spherical polynomial of degree $L \sim m^{1/(D-1)}$, constructed in the null space of the sampling operator. By the construction:
\[
  \|\nabla g\|_\infty \leq L \|g\|_\infty = L \cdot M L^{-(1+\alpha)} = M L^{-\alpha} = M m^{-\alpha/(D-1)} \to 0.
\]
Therefore the chord perturbation $|\delta C| \leq C M m^{-\alpha/(D-1)} \kappa R \to 0$ as $m \to \infty$.

For any noise level $\varepsilon > 0$: for $m$ large enough, $|\delta C| < \varepsilon$ at all $m$ X-ray directions. The bodies are X-ray-indistinguishable (within noise) but have $\delta_H \geq cMm^{-(1+\alpha)/(D-1)}$.

\textbf{Combined:} Both bounds give $\Theta(M m^{-(1+\alpha)/(D-1)})$. \qed
\end{proof}

\begin{remark}[Why the angular gap is the bottleneck, not the measurement type]
The upper bound uses only one fact about the measurement: it provides \emph{some} constraint near each $u_i$. Whether that constraint comes from a width (one scalar), an X-ray (a function), or a full tomographic scan (the complete density in a strip) is irrelevant. The bottleneck is the gap between measurement directions, where the reconstruction relies purely on regularity. No measurement at direction $u$ can help reconstruct $h$ at a direction $v$ far from $u$, beyond what the regularity of $h$ already guarantees.
\end{remark}
